% ============================================================
%  sections/intro.tex  –  Introduction
% ============================================================

Over the last two decades, Mexico has repeatedly revisited cannabis policy through criminal-law reform, public-health legislation, and constitutional adjudication.  The 2009 \emph{narcomenudeo} reform decriminalized small-quantity possession and reassigned retail drug offenses to state courts \citep{dof2009narcomenudeo,lgs2009doseTable}.  In 2017, amendments to the General Health Law authorized medical and scientific uses \citep{dof2017medicalcannabis}.  Most recently, the Supreme Court's 2021 General Declaration of Unconstitutionality (DGI~1/2018) invalidated the blanket prohibition on recreational use, mandating comprehensive regulation \citep{scjn2021dgi}.

These legal shifts have been accompanied by sustained attention to the fiscal design of a regulated market.  The Senate approved in 2020 the \emph{Ley Federal para la Regulaci\'{o}n del Cannabis}, whose fiscal chapter, building on an earlier legislative initiative \citep{delgado2019iniciativa}, envisioned a 12\% \emph{ad valorem} excise tax plus a per-gram surcharge under the IEPS framework \citep{senado2021minuta}.  The Chamber of Deputies passed an amended version in March~2021, but disagreements between the two chambers left the bill without definitive enactment \citep{diputados2021dictamen}.  Comparative analyses have evaluated international cannabis-tax regimes to inform domestic design \citep{clavellina2020cannabis_taxes}, academic work has projected the fiscal-revenue implications of legalization \citep{hernandezmacías2021impacto}, and microeconomic studies have examined the theoretical conditions for optimal taxation of marijuana in Mexico \citep{perezromero2014elementos}.  Designing any such tax regime, and more broadly anticipating the market consequences of regulation, requires an empirical input that, to date, does not exist for Mexico: the price elasticity of demand for marijuana.

The absence of domestic estimates reflects a fundamental data constraint.  In illicit markets, transaction-level prices and quantities are rarely observed at scale; administrative price series are either absent or unreliable \citep{caulkins1998price}, and household surveys such as Mexico's ENCODAT capture consumption prevalence but not the price--quantity margin \citep{encodat2025}.  The international literature offers meta-analytic consensus elasticities in the range of $-0.30$ to $-0.56$ \citep{gallet2014monkey}, complemented by crowdsourced estimates for the United States \citep{davis2016crowd}, regionally disaggregated analyses \citep{halcoussis2017geo}, online-survey evidence from South Africa \citep{riley2020southafrica}, and recent evidence from legal retail markets \citep{han2025commercialization,dodd2025canada}.  What remains missing is a Mexico-specific estimate based on realized transactions, in a market where policy reform, supply-chain geography, quality segmentation, and enforcement conditions differ from the settings studied to date.

This paper fills the gap by exploiting a rare data source: an anonymous survey, fielded between February~19 and September~15, 2019 and promoted mainly through \emph{Vice en Espa\~{n}ol}, that asked respondents to report completed marijuana purchases in Mexico \citep{bejaranoromero2020dataset}.  The raw file includes response identifiers and timestamps; respondent age, gender, education, and self-reported socioeconomic status; state and municipality of purchase; purchase date; perceived product quality (seedless or seeded); quality-specific quantity and expenditure fields; the buyer's relationship to the seller; age of first cannabis consumption; purchase frequency; and usual purchase amount.  For each retained record, we coalesce the seeded and seedless quantity/expenditure fields into one transaction-level quantity and value, then compute unit price as pesos per gram.  After cleaning, the analytic sample contains 3{,}293 purchases dated January~1--September~11, 2019, spanning all 32 states and 351 municipalities.  The survey is not nationally representative; its value is that it observes realized prices, quantities, quality, timing, and geography at a level of detail unavailable in conventional Mexican data sources.

We estimate log--log demand specifications that yield a stable transaction-level price elasticity of approximately $-0.65$ across OLS models with progressively richer controls. This elasticity describes how the quantity purchased in a completed transaction varies with the unit price recorded for that transaction; it does not directly identify participation responses, purchase frequency, or total individual consumption.  Quality segmentation reveals meaningful heterogeneity: seedless (\emph{good}) marijuana has an elasticity of about $-0.61$, whereas seeded (\emph{bad}) marijuana is more price responsive at roughly $-0.89$.  The benchmark estimate remains similar under extreme-value trimming, heaping checks, municipality fixed effects, wild cluster bootstrap inference, and coefficient-stability diagnostics.  We also implement instrumental-variables diagnostics in three families: OpenStreetMap (OSM) road-network driving time to major production states computed with the Open Source Routing Machine (OSRM); one-month-lag seizure-gravity measures based on OSM driving time and M\'{e}xico Unido contra la Delincuencia (MUCD) marijuana-seizure data, derived from official responses to public-information requests; and external Hausman price instruments that average prices in other markets within the same quality--month cell.  The OSM production-hub instrument remains too weak to serve as an alternative benchmark.  The hub-only lagged seizure-gravity specification reproduces the benchmark magnitude, while stricter outside-region seizure gravity illustrates the tradeoff between a cleaner exclusion story and weaker first-stage strength.  The external Hausman instruments are stronger and produce negative elasticities.  We interpret these IV results as diagnostics that triangulate the benchmark rather than as a definitive causal replacement for it.

The study contributes to the literature in three ways.  First, it provides the first Mexico-specific transaction-level estimate of the price elasticity of marijuana demand, directly informing a policy debate in which cannabis taxation and regulation have advanced faster than domestic demand evidence \citep{perezromero2014elementos,clavellina2020cannabis_taxes,senado2021minuta}.  Second, it shows that quality-tier heterogeneity is empirically important in this setting: seeded and seedless marijuana differ not only in prices but also in estimated demand responsiveness.  Third, it evaluates the promise and limits of crowdsourced transaction data for illicit-market demand estimation by pairing the survey with municipality-level geography, OSM/OSRM driving times, and MUCD enforcement records \citep{davis2016crowd,halcoussis2017geo,riley2020southafrica}.  The policy relevance is therefore practical but bounded.  An elasticity near $-0.65$ suggests that price-based regulation can affect purchase quantities without eliminating the tax base, while the stronger response for seeded marijuana warns that uniform taxes may affect quality tiers unevenly.  At the same time, the estimate comes from illicit-market transactions before national recreational regulation, so it should be treated as one empirical input into tax design rather than as a complete forecast of a legalized market.

The remainder of the paper is organized as follows: Section~\ref{sec:litreview} reviews the literature, Section~\ref{sec:data} describes the data and empirical strategy, Section~\ref{sec:results} presents the results, and Section~\ref{sec:conclusions} concludes with policy implications.
